\documentclass{article}
\usepackage[utf8]{inputenc}
\usepackage[T1]{fontenc}
\usepackage{geometry}
\usepackage{graphicx}
\usepackage{float}
\usepackage{hyperref}
\usepackage[english,indonesian]{babel}

\title{Optimasi Pemrosesan Template LaTeX Menggunakan Pendekatan Berbasis AI}
\author{Academic Support Agent}
\date{\today}

\begin{document}
\maketitle

\begin{abstract}

\end{abstract}

\chapter{PENDAHULUAN}


\chapter{TINJAUAN PUSTAKA}


% Original content (commented out):
% % This template has borrowed some code that has been used in both LACIAM and 
% International Colloquium on Mathematical Modeling in Epidemiology held at the School of Applied Mathematics of Getulio Vargas

\documentclass{article}
\usepackage[utf8]{inputenc}
\usepackage[a4paper, margin=2cm]{geometry}
\pagenumbering{gobble}

\usepackage{mathpazo}
\usepackage{graphicx} % Required for inserting images
\usepackage{blindtext}

\usepackage{amssymb}
\usepackage{authblk}
\usepackage{amsmath}

\usepackage{calc}
\usepackage{color}
\usepackage{amsfonts}
\usepackage{latexsym}
\usepackage{placeins}
\usepackage{hyperref}

% \addtolength{\voffset}{-3.5cm} \addtolength{\textheight}{4cm}

\renewcommand{\refname}{\Large{\textbf{Bibliografia}}}
\renewcommand\Authfont{\scshape\small}
\renewcommand\Affilfont{\itshape\small}
\newcommand{\palavraschave}[1]{\vspace{.2cm}\noindent{\large{\bf Palavras-chave:}} #1\\}
\newcommand{\keywords}[1]{\vspace{.2cm}\noindent{\large{\bf Keywords:}} #1\\}
\setlength{\affilsep}{1em}
\newcommand{\smalllineskip}{\baselineskip=15pt}
\newcommand{\emailaddress}[1]{{\sf#1}}
\newcommand{\speaker}[1]{\author{\underline{#1}}}
\newcommand{\speakeraffil}[1]{\affil{#1}}
\let\LaTeXtitle\title
\renewcommand{\title}[1]{\LaTeXtitle{\Large{\textbf{#1}}}}

\title{Título da apresentação} % Write the tile of your contribution here

\date{\small III Semana de Extensão da Escola de Matemática Aplicada da FGV - SEEMAp\\
18 de outubro de 2023} % Do not modify this line


%%AUTHORS 
\speaker{Crazy Frog} % Write speaker name here

\author[2]{Taylor Swift} % Write first co-author name here, if there is any. If not, remove the line.

%%AFFILIATIONS
\speakeraffil{Escola de Matemática Aplicada, Fundação Getulio Vargas, Rio de Janeiro,\,\emailaddress{crazy.frog@fgv.edu.br}} % Write affiliation/s of the speaker here.

\affil[2]{Instituto/Universidade, Cidade,\,\emailaddress{swift@music.co}} % Write affiliation/s of the first co-author here, if there is any. If not, remove the line.

\begin{document}

\maketitle

\section*{Resumo}

\blindtext % Put your abstract in Portuguese here

\palavraschave{Música, dança, vocais.}

\section*{Abstract}

\blindtext % Put your abstract in English here\\

\keywords{Music, dance, vocals.}
%% IF YOU WANT TO INCLUDE A BIBLIOGRAPHY LIST, PUT THE CITATIONS IN THE SEQUEL

\bibliographystyle{plain}

\begin{thebibliography}{99}

% FOR ARTICLE REFERENCE FOLLOW THIS FORMAT:
\bibitem{ref1} {\sc J. Timberlake,}  {\it Article title.} Journal title,
{Volume, Number, Pages, Year.}

% FOR BOOK REFERENCE FOLLOW THIS FORMAT:
\bibitem{ref2} {\sc J. Timberlake, M. Jackson,}  {\it Book title},
{ Editor, City, Year.}

\end{thebibliography}

\end{document}


\end{document}